\setcode{utf8}
\thispagestyle{plain}
\addcontentsline{toc}{chapter}{Résumé en arabe}
\markboth{Résumé en arabe}{Résumé en arabe}

\begin{center}
{\Huge\bfseries\RL{ملخص}}
\end{center}

\vspace{0.6cm}

\footnotesize
\setstretch{1.0}

\begin{arabtext}
تظل عملية تحديث التطبيقات البرمجية مكلفة من حيث الوقت والخبرة، لأنها تجمع بين فهم الأنظمة القائمة، وحل مشكلات توافق الإصدارات، وتحويل الشيفرة، وتشغيل الأدوات، والتحقق، وتكرار دورات التشخيص. يهدف مشروع نهاية الدراسة هذا، المنجز داخل شركة \LR{CGI Technologies et Solutions Maroc}، إلى تصميم منصة \LR{Migration Factory} ذكية وقابلة للتوسع ومحكومة، قادرة على أتمتة جزء مهم من عملية الترحيل مع الإبقاء على تحكم المطور في القرارات الحساسة.

تعتمد الحلول المقترحة على معمارية هجينة متعددة الوكلاء ومنظمة مركزيا. تُخصَّص الوكلاء للمهام التي تتطلب الاستدلال أو التكيف، ولا سيما التحليل، والتخطيط، والمراجعة، والتحويل المعتمد على الوكلاء، والإصلاح، والمساعدة. أما العمليات الحرجة والقابلة للتحقق فتبقى تحت مسؤولية خدمات حتمية تتحكم في حالة سير العمل، والانتقالات، والأوامر، وعمليات الكتابة، والتحقق، وحفظ الأدلة. وقد أُنجز المشروع بصورة تدريجية من خلال تحليل الحاجة، وتصميم المعمارية، وتنفيذ المكونات، والتحقق التقني، وتشغيل سيناريوهات ممثلة، ثم إدماج الملاحظات بصورة متتابعة.

تم تطوير تطبيقين مستقلين. يستهدف الأول \LR{Java/Spring Boot}، وهو التطبيق الأكثر نضجا، إذ يدعم مسارا تدريجيا للترحيل، وبيئات عمل معزولة، والتحقق بواسطة \LR{Maven}، وحلقة إصلاح محكومة، وواجهة متابعة، ومساعدا، وتقارير. كما يدمج \LR{Azure OpenAI} عبر إعدادات خارجية تربط عمليات نشر النماذج بأدوار الوكلاء. أما التطبيق الثاني فيستهدف \LR{Angular} بهدف اختبار قابلية نقل المعمارية إلى منظومة تقنية أخرى. وتثبت أدلة التنفيذ المعتمدة نجاح الانتقال من \LR{Angular 18} إلى \LR{Angular 19} ثم من \LR{19} إلى \LR{20}، بينما تم إعداد مرحلة \LR{20} إلى \LR{21} دون بدء تنفيذها بعد، وما زال المسار الكامل من \LR{Angular 11} إلى \LR{21} يحتاج إلى الاستكمال.

تُبرز عملية التقييم أساسا أهمية الفصل بين الاستدلال، واتخاذ القرار، والتنفيذ من أجل تحسين قابلية التتبع، وإعادة الإنتاج، وحوكمة عمليات الترحيل. ومع ذلك، لم تُثبت بعد بصورة كمية المكاسب الإجمالية في الوقت والتكلفة والإنتاجية على مجموعة ممثلة من المشاريع. لذلك تتمثل المساهمة الرئيسية للمشروع في معمارية محكومة ترفع مستوى الأتمتة مع الحفاظ على ضوابط حتمية، وأدلة تقنية، وإشراف بشري.

الكلمات المفتاحية: \LR{Angular}، تحديث التطبيقات، \LR{Azure OpenAI}، الإنسان ضمن الحلقة، \LR{Java}، النماذج اللغوية الكبيرة، الأنظمة متعددة الوكلاء، التنسيق، ترحيل البرمجيات، \LR{Spring Boot}، قابلية التتبع.
\end{arabtext}

\normalsize
\setstretch{1.12}
